\documentclass[a4paper,11pt]{article}
\usepackage[utf8]{inputenc}
\usepackage[T1]{fontenc}
\usepackage[brazil]{babel}
\usepackage{amsmath,amssymb}
\usepackage[margin=2.5cm]{geometry}
\usepackage[hidelinks]{hyperref}

\title{Equilibração e termalização:\\das cadeias de spins ao modelo de Hubbard}
\author{}
\date{19 de novembro de 2026}

\begin{document}
\maketitle

\section{Contexto e objetivo}

Este roteiro registra uma recomendação inicial de estudo para quem já tem
familiaridade com mecânica quântica e operadores e vem de mecânica estatística
e fenômenos críticos. A transição proposta é da descrição de estados de
equilíbrio para a pergunta dinâmica: como observáveis de um sistema quântico
preparado fora do equilíbrio podem adquirir valores aproximadamente
estacionários e, eventualmente, térmicos?

O foco provisório será em \textbf{sistemas fechados, com evolução unitária}.
Isso permite estudar o papel das interações, das conservações e da estrutura
espectral sem introduzir inicialmente um banho. Não é uma decisão
irreversível sobre o projeto: sistemas abertos poderão entrar se o objetivo
posterior exigir. A cadeia de spins será uma etapa intermediária para chegar
ao modelo de Hubbard, não uma substituição do objetivo original.

Adotaremos $\hbar=k_B=1$. O tempo será denotado por $\tau$; a letra $t$
ficará reservada ao hopping do modelo de Hubbard, sem fixar $t=1$.
Energias e tempos têm, portanto, unidades inversas entre si.

\section{Três noções que precisam permanecer distintas}

\begin{description}
  \item[Relaxação.] Aproximação aparente de um observável a um patamar,
  ou perda de um desvio inicial, dentro de uma janela temporal especificada.
  Uma curva com esse aspecto não demonstra, sozinha, equilibração ou
  termalização.
  \item[Equilibração.] Permanência do observável próximo de um valor de
  referência durante a maior parte dos tempos considerados, com flutuações
  pequenas segundo um critério declarado. Não exige convergência pontual
  quando $\tau\to\infty$.
  \item[Termalização.] Compatibilidade dos valores de observáveis locais
  com um ensemble térmico apropriado à energia e às conservações relevantes
  da preparação. É uma comparação adicional, não um sinônimo de estabilização.
\end{description}

Um estado global inicialmente puro permanece puro sob evolução unitária:
não precisa se transformar em uma matriz densidade de Gibbs para que
observáveis locais tenham comportamento térmico. Em sistemas finitos, a
discrição do espectro permite oscilações e recorrências; não se deve
prometer convergência exata nem inferir um limite termodinâmico de um
único tamanho.

\section{Por que não ficar apenas no dímero?}

O dímero fechado é útil para entender mecanismos locais, sinais e
frequências: poucos gaps de energia relevantes produzem oscilações e
recorrências facilmente identificáveis. Ele não constitui, porém, uma
plataforma robusta para investigar termalização genérica. Cadeias finitas
oferecem uma dinâmica mais rica e podem apresentar equilibração prática
em janelas acessíveis, mas exigem comparações entre tamanhos e tempos.
\textbf{Nem toda cadeia termaliza.}

\newpage
\section{Modelos e domínio da ponte entre spins e férmions}

\subsection{Cadeia de spins como ponto de partida}

Considere uma cadeia aberta com $L$ sítios, $L$ par, e spin $1/2$ em cada
sítio. Escrevemos $\mathbf S_i=\boldsymbol{\sigma}_i/2$, onde
$\boldsymbol{\sigma}_i$ reúne as matrizes de Pauli no sítio $i$. O Hamiltoniano
de partida será
\begin{equation}
  H_s =
  J_1\sum_{i=1}^{L-1}\mathbf S_i\cdot\mathbf S_{i+1}
  +J_2\sum_{i=1}^{L-2}\mathbf S_i\cdot\mathbf S_{i+2},
  \qquad J_1>0.
  \label{eq:spin}
\end{equation}
Os parâmetros $J_1$ e $J_2$ são acoplamentos de troca, com dimensão de
energia. O primeiro é antiferromagnético; o segundo será uma variável
de comparação, ainda sem valor escolhido.

Para $J_2=0$, temos a cadeia de Heisenberg integrável. Um $J_2$ genérico,
longe de pontos especiais, permite estudar uma versão não integrável.
Essa mudança não garante termalização, tampouco uma escala de relaxação
curta. Ambos os termos preservam a simetria global de spin
$\mathrm{SU}(2)$; em particular, conservam
$\mathbf S_{\mathrm{tot}}=\sum_i\mathbf S_i$.
Se estatísticas espectrais forem usadas mais adiante, será necessário
separar os setores de simetria pertinentes, inclusive as simetrias
espaciais presentes, em vez de misturar todos os níveis.

\subsection{O Hubbard tem outros graus de liberdade}

Uma cadeia de spins $1/2$ possui dois estados locais,
$|\uparrow\rangle$ e $|\downarrow\rangle$. No Hubbard de uma banda,
o espaço local tem quatro estados:
\[
  |0\rangle,\qquad |\uparrow\rangle,\qquad
  |\downarrow\rangle,\qquad |\uparrow\downarrow\rangle.
\]
Logo, o Hubbard contém dinâmica de carga e dupla ocupação que uma
cadeia apenas de spins não representa.

Para uma cadeia uniforme aberta com hopping entre primeiros vizinhos,
adotaremos o Hamiltoniano canônico
\begin{equation}
  H_{\mathrm H} =
  -t\sum_{\langle i,j\rangle,\sigma}
  \left(c^\dagger_{i\sigma}c_{j\sigma}
        +c^\dagger_{j\sigma}c_{i\sigma}\right)
  +U\sum_{i=1}^{L}n_{i\uparrow}n_{i\downarrow},
  \qquad U>0,
  \label{eq:hubbard}
\end{equation}
com $t$ real, $\sigma\in\{\uparrow,\downarrow\}$ e
$n_{i\sigma}=c^\dagger_{i\sigma}c_{i\sigma}$. Cada ligação não orientada
aparece uma única vez na soma; o conjugado hermitiano está explícito.
O número de partículas $N$ é fixo, a densidade é $n=N/L$ e meio
preenchimento significa $N=L$. Não incluímos o termo grand-canônico
$-\mu N$.

Em meio preenchimento, para $U\gg |t|$ e no subespaço apropriado de baixa
energia, com ocupação simples dominante, a projeção do Hubbard produz,
na ordem dominante, a troca antiferromagnética
\begin{equation}
  J=\frac{4t^2}{U}.
\end{equation}
Essa ponte vale para o hopping de primeiros vizinhos não frustrado aqui
considerado; não identifica os dois modelos para energias ou estados
iniciais arbitrários. Em particular, a cadeia com $J_2\ne0$ é uma
comparação dinâmica deliberada, não o modelo efetivo dominante desse
Hubbard de primeiros vizinhos.

\newpage
\section{Protocolo inicial de estudo}

\subsection{Preparação, evolução e observáveis}

Uma preparação simples é o estado produto de Néel,
\begin{equation}
  |\psi_0\rangle=|\uparrow\downarrow\uparrow\downarrow\cdots\rangle,
  \qquad
  |\psi(\tau)\rangle=e^{-iH_s\tau}|\psi_0\rangle.
\end{equation}
Para um operador $O$, escreveremos
$\langle O\rangle_\tau=\langle\psi(\tau)|O|\psi(\tau)\rangle$.
Dois observáveis iniciais são
\begin{align}
  m_s(\tau)&=\frac{1}{L}\sum_{i=1}^{L}
       (-1)^{i-1}\langle S_i^z\rangle_\tau,
       \label{eq:ms}\\
  C_i(\tau)&=\langle\mathbf S_i\cdot\mathbf S_{i+1}\rangle_\tau,
       \qquad 1\le i<L.
       \label{eq:corr}
\end{align}
A magnetização alternada mede a memória do padrão inicial; a correlação
$C_i$ acompanha uma ligação específica, cuja posição deve ser registrada
por causa das bordas abertas.

É importante não trocar automaticamente essa correlação local por uma
média espacial que possa ser conservada. Para $J_2=0$,
$J_1\sum_i C_i(\tau)=\langle H_s\rangle$ é constante.
A magnetização total $S_{\mathrm{tot}}^z$ também é conservada: serve
como controle de consistência, não como diagnóstico de relaxação.

\subsection{Separar estabilização de comparação térmica}

O percurso recomendado é:
\begin{enumerate}
  \item Observar as séries temporais de $m_s$ e de correlações por
  ligação, identificando transientes, oscilações e eventuais recorrências,
  sem nomear todo patamar como térmico.
  \item Comparar médias e flutuações em várias janelas temporais e em
  diferentes tamanhos $L$, mantendo o protocolo e a escolha de ligação
  comparáveis. Não fixamos aqui um tamanho nem uma janela.
  \item Repetir a análise para $J_2=0$ e para um $J_2$ genérico não
  integrável. Como o Hamiltoniano muda, a energia inicial e a referência
  de equilíbrio devem ser determinadas em cada caso.
  \item Só então comparar valores aparentemente equilibrados com
  previsões térmicas, explicitando a energia da preparação, os setores
  de simetria e as conservações usadas na construção do ensemble.
\end{enumerate}

Para tornar o segundo passo quantitativo, se $A(\tau)$ for uma das
médias quânticas acima, uma janela $W=[\tau_a,\tau_b]$ permite definir
\begin{equation}
  \overline A_W=\frac{1}{\Delta\tau}\int_{\tau_a}^{\tau_b}A(\tau)\,d\tau,
  \qquad
  \delta A_W=
  \left[\frac{1}{\Delta\tau}\int_{\tau_a}^{\tau_b}
  \bigl(A(\tau)-\overline A_W\bigr)^2\,d\tau\right]^{1/2},
  \quad \Delta\tau=\tau_b-\tau_a.
\end{equation}
Aqui $\delta A_W$ mede flutuações temporais da média quântica, não a
incerteza de uma medida individual. Sua escala de comparação e a
tolerância para chamar um patamar de estável ainda deverão ser escolhidas.

\newpage
\subsection{Uma referência não necessariamente térmica}

Em um sistema finito, a média temporal de tempo infinito da matriz
densidade tem a forma
\begin{equation}
  \overline\rho=\sum_E P_E\rho_0P_E,
  \qquad \rho_0=|\psi_0\rangle\langle\psi_0|,
  \label{eq:diagonal}
\end{equation}
onde $P_E$ projeta no autoespaço de energia $E$. A expressão preserva
coerências dentro de subespaços degenerados; não se deve descartá-las
por uma escolha arbitrária de base. Essa referência, usualmente
associada ao ensemble diagonal, não é automaticamente térmica.
Conhecê-la tampouco demonstra que as flutuações em uma janela acessível
sejam pequenas.

A integrabilidade traz conservações adicionais. Tanto a cadeia de
Heisenberg com $J_2=0$ quanto a cadeia uniforme de Hubbard com hopping
apenas entre primeiros vizinhos são integráveis. Por isso, quando
necessário, a referência estacionária deverá considerar essas
conservações, por exemplo mediante um ensemble de Gibbs generalizado
(GGE). Isso não significa que todo estado inicial ou todo observável
nesses modelos esteja impedido de apresentar valores compatíveis com
uma descrição térmica usual.

\section{Retorno ao Hubbard e critérios de progresso}

Ao retornar à Eq.~\eqref{eq:hubbard}, o objetivo será conservar as
distinções conceituais e acrescentar a dinâmica de carga. Além de
observáveis de spin, podem ser acompanhadas a dupla ocupação média
e, quando a preparação produzir um padrão alternado de densidade,
seu desequilíbrio:
\begin{equation}
  D(\tau)=\frac{1}{L}\sum_i
    \langle n_{i\uparrow}n_{i\downarrow}\rangle_\tau,
  \qquad
  I_n(\tau)=\frac{1}{L}\sum_i(-1)^{i-1}
    \langle n_{i\uparrow}+n_{i\downarrow}\rangle_\tau.
\end{equation}
No estado de Néel do Hubbard com uma partícula por sítio e $L$ par,
a densidade inicial é uniforme e $I_n(0)=0$. Assim, esse desequilíbrio
não deve ser anunciado como medida de relaxação de um padrão de carga
inicial nessa preparação. Um protocolo que o use com esse propósito
precisará preparar tal padrão.

O avanço nesta etapa será reconhecer o que cada observável informa,
separar conservação de relaxação e distinguir estabilidade temporal
de compatibilidade térmica. Evidências em tamanhos finitos deverão
continuar identificadas como tais, sem extrapolação automática.

As próximas decisões são escolher os parâmetros e o protocolo,
selecionar observáveis e referências de comparação e definir os
critérios temporais. Só depois será escolhido um método numérico em
Julia, precedido de estimativas da dimensão do espaço de Hilbert,
dos setores aproveitáveis, da memória e do custo. Não há compromisso
aqui com um valor numérico de $L$ ou com um algoritmo.

\paragraph{Escopo e pendências.}
Este arquivo é um roteiro de estudo, não uma revisão de literatura,
uma derivação extensa ou um relato de resultados. Nenhum cálculo
numérico ou teste foi executado para produzi-lo. O roteiro pessoal
do usuário ainda não foi disponibilizado e poderá ser incorporado
posteriormente, ajustando a sequência sem antecipar novas decisões
de escopo.

\end{document}
